\chapter{Achondroplasia}
\index{Achondroplasia}

\section*{Definition and Overview}
Achondroplasia is the most common form of disproportionate short stature (dwarfism). It is a genetic condition that affects the growth of bone and cartilage, producing short arms and legs, a relatively large head, and an average-sized trunk. The name means ``without cartilage formation'', although the underlying problem is really poor conversion of cartilage into bone within the growth plates, rather than an absence of cartilage.

Achondroplasia is present from birth and is usually diagnosed in childhood, or sometimes before birth, on the basis of characteristic physical features, X-ray findings, or genetic testing. People with the condition have normal intelligence, and with good medical care and management of complications, most lead full, active, and independent lives.

\section*{Epidemiology}
Achondroplasia is the most common skeletal dysplasia, occurring in roughly 1 in 20,000 to 30,000 births worldwide. It affects males and females equally and occurs across all ethnic groups. Around 80 per cent of cases arise from a new gene change (mutation) in families with no history of the condition, and the risk of a new mutation increases slightly with paternal age.

When a parent has achondroplasia, each child has a 1 in 2 chance of inheriting the condition. If both parents are affected, there is a 1 in 4 chance that a child inherits two copies of the changed gene, a severe form of the condition that is usually fatal in infancy.

\section*{Aetiology and Pathophysiology}
Achondroplasia is caused by a change in the \textit{FGFR3} gene, which carries the instructions for a protein that regulates bone growth. In about 80 per cent of cases the gene change is new (spontaneous), while in the remaining cases it is inherited from an affected parent.

The altered FGFR3 protein is overactive and slows the conversion of cartilage into bone within the growth plates of the long bones. As a result, the arms and legs grow much less than the trunk and skull, while the spine, trunk, and internal organs develop at a more normal pace. Because the gene change is present from conception, its effect on growth begins before birth and continues throughout childhood.

\section*{Clinical Features}
\begin{itemize}
    \item \textbf{Disproportionate short stature:} an average-sized trunk with short arms and legs
    \item \textbf{Distinctive facial features:} a relatively large head, prominent forehead, and flattened bridge of the nose
    \item \textbf{Hand and foot findings:} short fingers and toes, often with a gap between the middle and ring fingers (``trident'' hand)
    \item \textbf{Limb and spine alignment:} bowed legs and increased curvature of the lower back (lordosis)
    \item \textbf{Ear and breathing problems:} increased risk of middle-ear infections, hearing difficulties, snoring, and sleep apnoea
    \item \textbf{Motor development:} delayed motor milestones that are eventually reached, with normal intellectual development
\end{itemize}

\section*{Differential Diagnosis}
Several other conditions can cause short stature or a similar skeletal appearance.

\begin{itemize}
    \item Hypochondroplasia, a milder and related condition
    \item Thanatophoric dysplasia, a severe, usually fatal skeletal dysplasia
    \item Pseudoachondroplasia and other skeletal dysplasias
    \item Growth hormone deficiency or other hormonal causes of short stature
    \item Constitutional short stature in children of short parents
\end{itemize}

\section*{When to Seek Medical Care}
See a doctor if an infant or child with achondroplasia has breathing difficulty, poor feeding, or a delay in reaching developmental milestones. Medical review is needed for persistent back or leg pain, weakness, numbness, or difficulty walking, which can indicate narrowing of the spinal canal. Adults should seek review for any new neurological symptoms or signs of sleep disturbance such as loud snoring and daytime sleepiness.

\textbf{Contact your local emergency services for:}
\begin{itemize}
    \item Difficulty breathing, especially a baby or child turning blue or stopping breathing
    \item Sudden weakness or paralysis of the arms or legs
    \item Loss of consciousness or a seizure
    \item Sudden severe headache with vomiting, especially in an infant with a rapidly growing head
\end{itemize}

\section*{Diagnosis and Investigations}
\begin{itemize}
    \item \textbf{Clinical examination:} the pattern of short limbs with an average-sized trunk and typical facial features is usually recognisable
    \item \textbf{Measurement:} height, weight, and head circumference plotted on growth charts specific to achondroplasia
    \item \textbf{X-rays:} typical changes in the long bones, spine, and pelvis help confirm the diagnosis
    \item \textbf{Genetic testing:} a blood test that identifies the \textit{FGFR3} gene change confirms the diagnosis
    \item \textbf{Prenatal diagnosis:} sometimes suspected on ultrasound and confirmed by genetic testing during pregnancy
    \item \textbf{Specialist referral:} ongoing care by a team of paediatricians, orthopaedic surgeons, neurologists, and ear, nose, and throat specialists
\end{itemize}

\section*{Management}
\begin{itemize}
    \item \textbf{Multidisciplinary care:} regular monitoring of growth, development, hearing, breathing, and spinal health
    \item \textbf{Prompt treatment of ear infections:} to protect hearing and speech development
    \item \textbf{Sleep assessment:} investigation and treatment of obstructive sleep apnoea if present
    \item \textbf{Physiotherapy and occupational therapy:} to support movement, strength, and independence
    \item \textbf{Orthopaedic surgery:} for bowed legs, spinal deformity, or compression of the spinal cord
    \item \textbf{Limb-lengthening surgery:} an option some families choose, although it is lengthy, demanding, and controversial
    \item \textbf{Genetic counselling:} accurate information about inheritance and recurrence risks
\end{itemize}

\section*{Complications}
\begin{itemize}
    \item Narrowing around the spinal cord at the base of the skull (foramen magnum stenosis), which can compress the brainstem in infancy
    \item Spinal stenosis later in life, causing back pain, numbness, or weakness in the legs
    \item Hydrocephalus (a buildup of fluid in the brain), requiring monitoring
    \item Recurrent middle-ear infections with the risk of hearing loss
    \item Obstructive sleep apnoea
    \item Obesity, which is more common and adds to physical strain
    \item Bowed legs and early arthritis of the hips and knees
\end{itemize}

\section*{Prognosis and Outcomes}
Most people with achondroplasia have normal intelligence and a normal or near-normal life expectancy. The main complications, such as sleep apnoea, ear infections, and spinal problems, are treatable when detected early, so regular specialist follow-up is important across the lifespan.

The principal health risks occur in infancy, when narrowing around the brainstem and hydrocephalus can threaten life if not recognised. In adult life, attention to weight, back health, joint symptoms, and cardiovascular risk allows most people to work, drive, and live independently. Psychosocial support and community participation are important for overall wellbeing.

\section*{Prevention and Self-Care}
\begin{itemize}
    \item Attend all scheduled checks of growth, hearing, vision, and sleep in infancy and childhood
    \item Report any new back or leg pain, weakness, or numbness promptly
    \item Maintain a healthy weight through diet and exercise suitable for the condition, such as swimming
    \item Protect the head and neck, using appropriately sized car seats and boosters for children
    \item Adapt the home and workplace, including work-surface heights and reach aids, to reduce strain
    \item Seek genetic counselling before planning a pregnancy
    \item Encourage social participation and seek support for emotional wellbeing
\end{itemize}

\section*{Sources and Further Reading}
The following organisations provide reliable information on achondroplasia for families, patients, and health professionals.

\begin{itemize}
    \item NHS. \textit{Information on achondroplasia and short stature.} https://www.nhs.uk/conditions/
    \item Mayo Clinic. \textit{Overview of achondroplasia, its complications, and management.} https://www.mayoclinic.org/
    \item MedlinePlus. \textit{Health information on achondroplasia and genetic conditions.} https://medlineplus.gov/
    \item World Health Organization. \textit{Resources on congenital conditions and child health.} https://www.who.int/
    \item MSD Manual. \textit{Professional information on skeletal dysplasias.} https://www.msdmanuals.com/
\end{itemize}
